% !TEX program = pdflatex
\documentclass[11pt,a4paper]{article}

% --- Encoding & Vietnamese support (works out-of-the-box on Overleaf & TeX Live) ---
\usepackage[utf8]{inputenc}
\usepackage[T5]{fontenc}
\usepackage[vietnamese]{babel}

% --- Clean Sans-Serif Executive Typography ---
\renewcommand{\familydefault}{\sfdefault}

% --- Page Geometry & Spacing ---
\usepackage[a4paper, top=2.2cm, bottom=2.1cm, left=2.0cm, right=2.0cm, headheight=24pt]{geometry}
\usepackage{parskip}
\setlength{\parindent}{0pt}
\setlength{\parskip}{6pt}

% --- Colors & Styling (Pitch-Deck Palette) ---
\usepackage[table]{xcolor}
\definecolor{NavyDark}{HTML}{16425B}
\definecolor{TealAccent}{HTML}{2A8C96}
\definecolor{AmberAccent}{HTML}{E07A16}
\definecolor{SlateText}{HTML}{2B3440}
\definecolor{MutedGray}{HTML}{5C6B73}
\definecolor{LightBg}{HTML}{F4F7F9}
\definecolor{WarmBg}{HTML}{FFF8F0}
\definecolor{BorderGray}{HTML}{D6E0E7}

\color{SlateText}

% --- Graphics, Tables & Boxes ---
\usepackage{graphicx}
\graphicspath{{../figures/}{figures/}{./}}
\usepackage{float}
\usepackage[section]{placeins}
\usepackage{booktabs}
\usepackage{tabularx}
\usepackage{array}
\usepackage{tcolorbox}
\tcbuselibrary{skins,breakable}
\usepackage{enumitem}
\setlist[itemize]{leftmargin=1.4em, itemsep=3pt, topsep=2pt}

% --- Headers, Footers & Section Titles ---
\usepackage{fancyhdr}
\pagestyle{fancy}
\fancyhf{}
\renewcommand{\headrulewidth}{0.4pt}
\renewcommand{\headrule}{\hbox to\headwidth{\color{BorderGray}\leaders\hrule height \headrulewidth\hfill}}
\lhead{\footnotesize\color{MutedGray}\textbf{LẬP TRÌNH XỬ LÝ DỮ LIỆU \textbar{} NHÓM 07}}
\rhead{\footnotesize\color{TealAccent}\textbf{TORONTO AIRBNB}}
\lfoot{\footnotesize\color{MutedGray}Báo cáo Phân tích Thị trường Airbnb Toronto · Mẫu phân công · 20/11/2026}
\rfoot{\footnotesize\color{MutedGray}\thepage}

\usepackage{titlesec}
\titleformat{\section}
  {\Large\bfseries\color{NavyDark}}
  {\colorbox{TealAccent}{\color{white}\small\ \thesection\ }\hspace{0.5em}}
  {0pt}{}
\titlespacing*{\section}{0pt}{14pt}{6pt}

\titleformat{\subsection}
  {\normalsize\bfseries\color{TealAccent}}
  {}{0pt}{}
\titlespacing*{\subsection}{0pt}{10pt}{4pt}

% --- Hyperlinks ---
\usepackage[hidelinks]{hyperref}
\hypersetup{
  pdftitle={Báo cáo Phân tích Thị trường Airbnb Toronto},
  pdfauthor={Nhóm 07},
  pdfsubject={Nguồn cung, mặt bằng giá, lịch phòng và trải nghiệm khách}
}

% --- Custom Pitch-Deck Components ---
\newtcolorbox{storybox}[1][]{
  enhanced, breakable,
  colback=LightBg, colframe=TealAccent,
  boxrule=0pt, leftrule=3.5pt, arc=2pt,
  left=10pt, right=10pt, top=7pt, bottom=7pt,
  fonttitle=\bfseries\color{NavyDark},
  title={#1}
}

\newtcolorbox{caveatbox}[1][Lưu ý diễn giải]{
  enhanced, breakable,
  colback=WarmBg, colframe=AmberAccent,
  boxrule=0pt, leftrule=3.5pt, arc=2pt,
  left=10pt, right=10pt, top=7pt, bottom=7pt,
  fonttitle=\bfseries\color{AmberAccent},
  title={#1}
}

\newcommand{\kpicard}[3]{%
  \begin{tcolorbox}[
    colback=LightBg, colframe=BorderGray, boxrule=0.6pt, arc=3pt,
    left=6pt, right=6pt, top=8pt, bottom=8pt, halign=center
  ]
    {\LARGE\bfseries\color{NavyDark} #1}\\[3pt]
    {\footnotesize\bfseries\color{TealAccent} #2}\\[2pt]
    {\scriptsize\color{MutedGray} #3}
  \end{tcolorbox}%
}

% PDF-only figure macro. [H] keeps each piece of evidence beside its conclusion.
\newcommand{\analystfigure}[3][0.92\textwidth]{%
  \begin{figure}[H]
    \centering
    \IfFileExists{../figures/#2.pdf}{%
      \includegraphics[width=#1]{../figures/#2.pdf}%
    }{%
      \IfFileExists{figures/#2.pdf}{%
        \includegraphics[width=#1]{figures/#2.pdf}%
      }{%
        \IfFileExists{#2.pdf}{%
          \includegraphics[width=#1]{#2.pdf}%
        }{%
          \fbox{\parbox{0.86\textwidth}{\centering\vspace{1.2cm}\color{MutedGray}\small
            \textbf{[Thiếu biểu đồ PDF: \texttt{\detokenize{#2}.pdf}]}\\[4pt]
            Đặt thư mục \texttt{figures/} cùng cấp hoặc ở thư mục cha để tự động chèn biểu đồ.%
            \vspace{1.2cm}}}%
        }%
      }%
    }%
    \caption{\small\color{MutedGray}\textit{#3}}
  \end{figure}
}

\begin{document}
\thispagestyle{empty}

% =====================================================================
% HERO BANNER
% =====================================================================
\begin{tcolorbox}[
  colback=NavyDark, colframe=NavyDark, arc=4pt,
  left=16pt, right=16pt, top=16pt, bottom=16pt
]
  {\footnotesize\color{BorderGray}\textbf{BÁO CÁO PHÂN TÍCH THỊ TRƯỜNG \textbar{} NHÓM 07}}\\[6pt]
  {\raggedright\Huge\bfseries\color{white} Thị trường Airbnb Toronto:\\[2pt] Nguồn cung, giá \& trải nghiệm khách\par}\vspace{10pt}
  {\small\color{BorderGray}
    \textbf{Người đọc:} Cơ quan quản lý du lịch thành phố Toronto \quad\textbar\quad
    \textbf{Thực hiện:} Nhóm 07\\
    \textbf{Nguồn dữ liệu:} Inside Airbnb (CC BY 4.0) · 4 mốc chụp (2025-09 đến 2026-09)
  }
\end{tcolorbox}

\vspace{4pt}

% --- 3-COLUMN HEADLINE KPI STRIP ---
\noindent
\begin{minipage}[t]{0.32\textwidth}
  \kpicard{86.75\%}{CHỖ Ở NGUYÊN CĂN}{Mốc 2026-09-16 (n = 928)}
\end{minipage}\hfill
\begin{minipage}[t]{0.32\textwidth}
  \kpicard{201.07 CAD}{GIÁ TRUNG VỊ / ĐÊM}{Sau lọc QA (n = 783)}
\end{minipage}\hfill
\begin{minipage}[t]{0.32\textwidth}
  \kpicard{5.0 lần}{TẬP TRUNG NGUỒN CUNG}{Waterfront (161) vs Niagara (32)}
\end{minipage}

\vspace{6pt}

% --- EXECUTIVE SUMMARY ---
\begin{storybox}[Ý chính]
Tại mốc gần nhất (\textbf{2026-09-16}), \textbf{Waterfront Communities-The Island có 161 chỗ ở có giá hợp lệ, gấp 5.0 lần Niagara} ($n = 32$); đây là mức tập trung trong mẫu phân công, không phải mật độ theo diện tích hay ước lượng cho toàn bộ Toronto.

\medskip
Báo cáo này phục vụ \textbf{Cơ quan quản lý du lịch thành phố Toronto} với 4 đầu mục cụ thể: (1) xác định điểm tập trung nguồn cung và mức độ chênh lệch giá giữa các quận, (2) theo dõi biến động lịch phòng qua 4 mốc chụp, (3) tóm tắt phản hồi của khách kèm đối chứng hiệu năng LLM so với mô hình từ điển, và (4) rà soát các lỗi diễn giải dễ gặp khi đọc số liệu.
\end{storybox}

\begin{caveatbox}[Giới hạn dữ liệu cần lưu ý]
Dữ liệu lấy từ \textbf{Inside Airbnb} (CC BY 4.0; không thu thập trực tiếp từ \texttt{airbnb.com}) trên \textbf{mẫu được giảng viên phân công} qua 4 mốc chụp (\texttt{2025-09-11}, \texttt{2026-03-19}, \texttt{2026-06-15}, \texttt{2026-09-16}), nên \textbf{không đại diện cho toàn bộ thị trường thành phố}. Mốc cuối (\texttt{2026-09-16}) giảm \textbf{72 chỗ ở} và \textbf{235 đánh giá} so với 3 mốc đầu (chi tiết tại \texttt{missing\_ids.csv}) do ngừng niêm yết. Vì vậy, so sánh giữa các mốc phản ánh mẫu quan sát tại từng thời điểm, không phải một nhóm cố định (cohort).
\end{caveatbox}

% =====================================================================
% SECTION 1: DATA FOUNDATION
% =====================================================================
\section{Ba mốc đầu đủ 1,000 chỗ ở; mốc cuối còn 928}

Bộ dữ liệu gồm 3 bảng chính (\texttt{listings}, \texttt{calendar}, \texttt{reviews}) được thu thập ở 4 mốc chụp trong vòng một năm. Sau bước kiểm soát chất lượng (QA), quy mô từng bảng như sau:

\begin{table}[H]
  \centering
  \small
  \renewcommand{\arraystretch}{1.25}
  \begin{tabularx}{0.92\textwidth}{Xrrr}
    \toprule
    \rowcolor{NavyDark}
    \color{white}\textbf{Mốc chụp (Snapshot)} &
    \color{white}\textbf{Listing sau QA} &
    \color{white}\textbf{Calendar sau QA} &
    \color{white}\textbf{Review sau QA} \\
    \midrule
    2025-09-11 & 1,000 & 365,000 & 5,000 \\
    \rowcolor{LightBg}
    2026-03-19 & 1,000 & 365,000 & 5,000 \\
    2026-06-15 & 1,000 & 365,000 & 5,000 \\
    \rowcolor{LightBg}
    2026-09-16 & 928 & 338,720 & 4,765 \\
    \bottomrule
  \end{tabularx}
  \caption{\small Quy mô mẫu dữ liệu sau QA qua 4 mốc chụp.}
\end{table}

Ba mốc đầu giữ nguyên đủ \textbf{1,000 chỗ ở} và \textbf{5,000 đánh giá}; mốc cuối còn \textbf{928 chỗ ở} và \textbf{4,765 đánh giá}. Đường dẫn tải và mã băm SHA256 của toàn bộ 12 tệp nguồn được lưu tại \texttt{download\_manifest.csv}.

% =====================================================================
% SECTION 2: SPATIAL & PRICING STORY
% =====================================================================
\section{Giá trung vị giữa 12 khu đông nhất chênh 2.6 lần}

\begin{storybox}[Kết luận chính]
Giá và nguồn cung phân hóa mạnh: trong 12 khu đông nhất, trung vị giá trải rộng từ \textbf{128.61 CAD/đêm} (Dovercourt-Wallace Emerson-Junction) đến \textbf{336.12 CAD/đêm} (Niagara), tức gấp \textbf{2.6 lần}. Tuy nhiên, chỗ ở nguyên căn (\textbf{232.88 CAD/đêm}) cũng đắt gấp \textbf{2.6 lần} phòng riêng (\textbf{90.5 CAD/đêm}). Do đó, khi quy hoạch hoặc đánh giá mức độ đắt đỏ theo quận, cơ quan quản lý cần tách riêng yếu tố \textbf{vị trí} khỏi \textbf{tỷ trọng nguyên căn} --- đồng thời lưu ý rằng 76/118 khu có $n < 5$ nên không đủ cỡ mẫu để xếp hạng đại diện.
\end{storybox}

\begin{itemize}
  \item \textbf{Chênh lệch giá theo khu vực (K1):} Trong 12 khu vực hành chính đông nhất tại mốc \texttt{2026-09-16}, \textbf{Niagara} có giá trung vị cao nhất (\textbf{336.12 CAD/đêm}, $n = 32$), gấp \textbf{2.6 lần} khu thấp nhất nhóm là \textbf{Dovercourt-Wallace Emerson-Junction} (\textbf{128.61 CAD/đêm}, $n = 24$).
  \item \textbf{Ba khu vực tập trung nguồn cung nhiều nhất (có giá hợp lệ):}
    (1) \textbf{Waterfront Communities-The Island} ($n = 161$, trung vị \textbf{279.17 CAD/đêm});
    (2) \textbf{Niagara} ($n = 32$, trung vị \textbf{336.12 CAD/đêm});
    (3) \textbf{Annex} ($n = 30$, trung vị \textbf{321.67 CAD/đêm}).
  \item \textbf{Ảnh hưởng của cơ cấu loại phòng:} Tại mốc cuối, trung vị giá của loại hình nguyên căn đạt \textbf{232.88 CAD/đêm}, gấp \textbf{2.6 lần} phòng riêng (\textbf{90.5 CAD/đêm}). Vì cơ cấu loại phòng không đồng đều giữa các quận, chênh lệch giá trung vị giữa các khu vực phản ánh đồng thời cả vị trí lẫn tỷ trọng nguyên căn, không nên quy kết đơn thuần do vị trí.
  \item \textbf{Hai cấp độ không gian phân tích:}
    (1) Các KPI từ bảng listing/lịch (K1, K2, K4, K6, K8) nhóm theo khu hành chính \texttt{neighbourhood\_cleansed} --- có \textbf{118} khu có giá ở mốc cuối nhưng \textbf{76} khu có $n < 5$, nên biểu đồ chỉ xếp hạng 12 khu đông nhất;
    (2) Bản đồ phân bố và K7 (LLM) nhóm theo 8 ô phân tích (\texttt{analysis\_area} A01--A08, mỗi vùng 110--130 listing) để đảm bảo đủ cỡ mẫu văn bản.
\end{itemize}

\analystfigure{fig_space_02_area_price}{Giá đêm trung vị (K1) tại 12 khu vực có cỡ mẫu lớn nhất, mốc 2026-09-16.}

\analystfigure{fig_space_01_listings_map}{Phân bố không gian theo trung vị giá (201.07 CAD/đêm) và 8 vùng phân tích A01--A08.}

% =====================================================================
% SECTION 3: TEMPORAL & CALENDAR DYNAMICS
% =====================================================================
\section{Ngày không khả dụng ổn định quanh 41.2\%--43.6\%}

\begin{storybox}[Kết luận chính]
Tỷ lệ ngày không khả dụng toàn mẫu đi ngang ở mức \textbf{41.18\%--43.62\%} qua cả 4 mốc, nhưng luôn vọt lên \textbf{57.64\%--71.34\%} ở tháng ngay sát ngày thu thập trước khi hạ nhiệt ở các tháng xa hơn và nhích lại vào mùa hè--thu. Cùng giai đoạn, nhóm thuê dài hạn $\ge 30$ đêm (K4) ổn định quanh \textbf{9.6\%--10.4\%} còn tỷ lệ chỗ ở chưa có review (K6) giảm từ \textbf{5.0\%} xuống \textbf{2.3\%}. Các chỉ số này hữu ích để theo dõi áp lực lịch phòng ngắn hạn, nhưng không thể dùng trực tiếp để ước tính công suất thuê hay doanh thu.
\end{storybox}

\begin{itemize}
  \item \textbf{Xu hướng khả dụng (K2, K3):} Tỷ lệ ngày không khả dụng toàn mẫu nằm gọn trong khoảng \textbf{41.18\%--43.62\%} qua 4 mốc chụp. Khi tách theo từng tháng của lịch 365 ngày tới (K3), tỷ lệ này luôn cao nhất ở tháng sát mốc chụp (\textbf{57.64\%--71.34\%}), giảm dần ở các tháng xa rồi tăng nhẹ trở lại vào cao điểm hè--thu.
  \item \textbf{Các chỉ số đi kèm (K4, K5, K6):} Tỷ lệ yêu cầu thuê tối thiểu $\ge 30$ đêm (K4) giữ quanh mức 9.6\%--10.4\%; tỷ lệ listing chưa có đánh giá (K6) giảm từ 5.0\% xuống 2.3\% ở mốc cuối. Tốc độ đánh giá trong mẫu (K5) dao động theo quý và đạt đỉnh \textbf{0.42 review/listing/quý tại 2025Q3}; đây không phải số lượt thuê hay toàn bộ review của Toronto.
\end{itemize}

\begin{caveatbox}[Giới hạn diễn giải bắt buộc]
(1) Ngày \texttt{available = f} gộp chung cả ngày khách đã đặt lẫn ngày chủ nhà tự khóa lịch, nên \textbf{không được đồng nhất với công suất thuê phòng thực tế hay doanh thu}.\\
(2) Số đánh giá (K5) không phải tổng lượt khách thuê, chỉ tính trên mẫu review được giao; tháng biên và quý cuối chưa đủ ngày, đồng thời thiếu dữ liệu từ các chỗ ở đã ngừng niêm yết.
\end{caveatbox}

\analystfigure{fig_time_01_occupancy_by_snapshot}{Tỷ lệ ngày không khả dụng theo tháng lịch (K3) qua 4 mốc chụp.}

\analystfigure{fig_time_02_kpi_trend}{Xu hướng các chỉ số K2, K4 và K6 qua 4 mốc chụp (tách khung chung trục thời gian).}

% =====================================================================
% SECTION 4: GUEST VOICE & AI BENCHMARK
% =====================================================================
\section{Baseline thắng cảm xúc; LLM mạnh hơn ở khía cạnh khen}

\begin{storybox}[Kết luận chính]
Trải nghiệm lưu trú nhìn chung rất tích cực (\textbf{89.7\%--95.3\%} ở cả 8 vùng phân tích), tập trung khen ở 3 yếu tố: \textbf{chủ nhà}, \textbf{vị trí} và \textbf{tiện nghi}. Về mặt công cụ, kiểm chứng trên 100 nhãn tay cho thấy không nên dùng một mô hình duy nhất cho mọi tác vụ: \textbf{Baseline từ điển} thắng rõ ở phân loại cảm xúc tổng thể (Accuracy \textbf{0.99 vs 0.91}), trong khi \textbf{LLM} vượt trội ở trích xuất khía cạnh khen (Micro-F1 \textbf{0.8503 vs 0.6934}). Dự án dùng \textbf{Gemini API free tier}, nên chi phí API thực tế cho mẫu 1,000 review là \textbf{0 USD}.
\end{storybox}

Toàn bộ \textbf{1,000/1,000} đánh giá hợp lệ được xử lý bằng \texttt{gemini-3.5-flash-lite} (prompt \texttt{v1}, \texttt{temperature=0}, lưu đệm JSONL) và đối chứng trên \textbf{100 nhãn gán tay} (\texttt{hand\_labels.csv}):

\begin{table}[H]
  \centering
  \small
  \renewcommand{\arraystretch}{1.25}
  \begin{tabularx}{0.92\textwidth}{Xlrrr}
    \toprule
    \rowcolor{NavyDark}
    \color{white}\textbf{Phương pháp} &
    \color{white}\textbf{Thước đo} &
    \color{white}\textbf{Giá trị} &
    \color{white}\textbf{Khoảng tin cậy 95\%} &
    \color{white}\textbf{Cỡ mẫu ($n$)} \\
    \midrule
    LLM & \texttt{accuracy} & 0.9100 & {[0.8500, 0.9600]} & 100 \\
    \rowcolor{LightBg}
    LLM & \texttt{macro\_f1} & 0.7112 & {[0.3139, 0.8070]} & 100 \\
    BASELINE & \texttt{accuracy} & 0.9900 & {[0.9700, 1.0000]} & 100 \\
    \rowcolor{LightBg}
    BASELINE & \texttt{macro\_f1} & 0.9679 & {[0.5799, 1.0000]} & 100 \\
    \bottomrule
  \end{tabularx}
  \caption{\small Đối chứng Accuracy và Macro-F1 giữa LLM và Baseline từ điển trên 100 nhãn tay (bootstrap 95\%).}
\end{table}

\analystfigure{fig_llm_03_model_vs_baseline}{Đối chứng Accuracy và Macro-F1 cảm xúc giữa LLM và Baseline từ điển (bootstrap 95\%).}

\begin{itemize}
  \item \textbf{Phân loại cảm xúc (\texttt{sentiment}):} Baseline đạt Accuracy cao hơn LLM có ý nghĩa thống kê (chênh lệch ghép cặp \texttt{llm\_minus\_baseline} = \textbf{-0.0800}, 95\% CI \texttt{[-0.1400, -0.0300]}; LLM lệch \textbf{9 ca}, Baseline lệch \textbf{1 ca}). Cả 9 ca lệch của LLM đều nằm ở ranh giới \texttt{mixed\_confusion} giữa \texttt{positive} và \texttt{mixed}, khi khách khen tổng thể nhưng thêm một lưu ý nhẹ.
  \item \textbf{Trích xuất khía cạnh khen và chê:} Ở khía cạnh khen (\texttt{aspects\_praise}), LLM đạt Micro-F1 = \textbf{0.8503} so với \textbf{0.6934} của Baseline. Ngược lại, ở khía cạnh bị chê (\texttt{aspects\_complaint}), do tập nhãn tay chỉ có \textbf{7/100} review có lời chê và LLM hay gắn nhãn chê cho các lưu ý thông tin nhẹ, Micro-F1 của LLM chỉ đạt \textbf{0.3636} so với \textbf{0.5333} của Baseline.
  \item \textbf{Giới hạn cỡ mẫu nhãn tay:} Tập đối chứng 100 review lệch lớp mạnh (\textbf{93} \texttt{positive}, \textbf{6} \texttt{mixed}, \textbf{1} \texttt{negative}), kéo khoảng tin cậy 95\% của Macro-F1 giãn rộng (\texttt{[0.3139, 0.8070]} với LLM); cần bổ sung thêm mẫu \texttt{negative}/\texttt{mixed} trước khi kết luận khái quát hơn.
  \item \textbf{Chi phí API thực tế:} Nhóm chỉ xử lý \textbf{1,000 review} trong mẫu LLM bằng \textbf{Gemini API free tier}, nên không phát sinh thanh toán (\textbf{0 USD}).
\end{itemize}

% =====================================================================
% SECTION 5: ANALYST AUDIT
% =====================================================================
\section{Cả 5 nhận định AI đều cần sửa hoặc bổ sung căn cứ}

\begin{storybox}[Kết luận chính]
Trong 5 nhận định phổ biến dễ gặp khi đọc nhanh báo cáo, \textbf{không có nhận định nào đạt yêu cầu giữ nguyên} (\textbf{4 nhận định phải sửa}, \textbf{1 nhận định chưa đủ căn cứ}). Điểm chung của các sai sót không nằm ở phép tính sai, mà ở việc \textbf{nhảy cóc từ số liệu mô tả sang kết luận nhân quả hoặc doanh thu} (như đồng nhất \texttt{available=f} với lượt đặt phòng, hoặc dùng giá trung bình 289.01 CAD bị kéo lệch bởi ngoại lai).
\end{storybox}

Nhóm áp dụng quy trình thẩm định 4 bước (Bài 14) --- \textbf{(1) Tính lại số liệu $\rightarrow$ (2) Kiểm tra phương pháp $\rightarrow$ (3) Kiểm tra diễn giải $\rightarrow$ (4) Kết luận: Giữ / Sửa / Chưa đủ căn cứ} --- cho 5 nhận định sau:

\begin{table}[H]
  \centering
  \footnotesize
  \renewcommand{\arraystretch}{1.28}
  \begin{tabularx}{\textwidth}{>{\raggedright\arraybackslash}p{3.6cm} l >{\raggedright\arraybackslash}p{2.3cm} >{\raggedright\arraybackslash}p{2.6cm} >{\raggedright\arraybackslash}X}
    \toprule
    \rowcolor{NavyDark}
    \color{white}\textbf{Nhận định cần thẩm định} &
    \color{white}\textbf{Số liệu} &
    \color{white}\textbf{Phương pháp} &
    \color{white}\textbf{Diễn giải} &
    \color{white}\textbf{Kết luận \& Bản sửa đúng mức} \\
    \midrule
    1. 86.75\% chỗ ở là nguyên căn, chứng tỏ tiềm năng dòng tiền đầu tư lớn. &
    Đúng & Đúng &
    Sai: suy diễn doanh thu/dòng tiền từ số lượng listing. &
    \textbf{Sửa:} Mốc 2026-09-16 có 928 chỗ ở trong mẫu, trong đó 86.75\% là nguyên căn (không có dữ liệu doanh thu). \\
    \rowcolor{LightBg}
    2. Giá trung bình (289.01 CAD) cao hơn trung vị (201.07 CAD) chứng minh thị trường siêu sang. &
    Đúng &
    Lệch: trung bình bị kéo bởi vài ngoại lai cực đoan. &
    Quá mức: dùng từ ``chứng minh''. &
    \textbf{Sửa:} Phân phối giá lệch phải; báo cáo dùng trung vị sau lọc IQR (\textbf{201.07 CAD/đêm}, $n = 783$) để phản ánh mức giá điển hình. \\
    3. Tỷ lệ không khả dụng 41.18\%--43.62\% cho thấy công suất lấp đầy khách luôn vượt 41\%. &
    Đúng &
    Sai: gộp ngày chủ nhà khóa lịch với ngày đã đặt. &
    Sai: đồng nhất \texttt{available=f} với lượt đặt phòng. &
    \textbf{Chưa đủ căn cứ} kết luận công suất thuê; chỉ ghi nhận tỷ lệ ngày không khả dụng trên lịch. \\
    \rowcolor{LightBg}
    4. Dovercourt-Wallace Emerson-Junction rẻ hơn Niagara là do kém hấp dẫn về vị trí. &
    Đúng &
    Thiếu: chưa kiểm soát cơ cấu loại phòng; 76/118 khu có $n < 5$. &
    Sai: suy diễn nhân quả đơn biến. &
    \textbf{Sửa:} Trong 12 khu đông nhất, trung vị giá dao động từ 128.61 đến 336.12 CAD/đêm và chịu ảnh hưởng của tỷ trọng nguyên căn. \\
    5. LLM vượt trội hoàn toàn mô hình từ điển ở mọi tác vụ phân tích đánh giá. &
    Sai ở \texttt{sentiment} &
    Đúng (bootstrap ghép cặp) &
    Sai: bỏ qua độ lệch lớp và kết quả thua trên cảm xúc. &
    \textbf{Sửa:} Baseline thắng có ý nghĩa thống kê trên cảm xúc tổng thể (Accuracy 0.99 vs 0.91); LLM chỉ vượt trội ở trích xuất khía cạnh khen (Micro-F1 0.85 vs 0.69). \\
    \bottomrule
  \end{tabularx}
  \caption{\small Bảng thẩm định 5 nhận định phân tích theo quy trình 4 bước.}
\end{table}

% =====================================================================
% SECTION 6: RESOLUTION
% =====================================================================
\section{Kết luận: ưu tiên nơi tập trung nguồn cung, không suy diễn công suất}

\begin{storybox}[Hàm ý cho cơ quan quản lý]
Trong mẫu phân công, \textbf{Waterfront Communities-The Island là khu cần được ưu tiên theo dõi về mức tập trung nguồn cung} (161 chỗ ở có giá hợp lệ, gấp 5.0 lần Niagara). Tuy nhiên, chênh lệch giá giữa các khu phải được đọc cùng cơ cấu loại phòng và cỡ mẫu; tỷ lệ \texttt{available=f} chỉ phản ánh ngày không khả dụng, \textbf{không đủ căn cứ để kết luận công suất thuê hoặc doanh thu}.
\end{storybox}

\begin{itemize}
  \item \textbf{Giữ để ra quyết định:} số chỗ ở theo khu, trung vị giá trên mẫu hợp lệ, tỷ lệ loại phòng và cỡ mẫu đi kèm.
  \item \textbf{Theo dõi nhưng không suy diễn:} ngày không khả dụng và số review là tín hiệu mô tả, không phải lượt đặt phòng hay số khách.
  \item \textbf{Cần bổ sung trước khi khái quát:} dữ liệu toàn thành phố, diện tích khu để tính mật độ, trạng thái đặt/khóa lịch và thêm nhãn \texttt{negative}/\texttt{mixed} cho đánh giá LLM.
\end{itemize}

% =====================================================================
% SECTION 7: METHODOLOGY, QA & REPRODUCIBILITY APPENDIX
% =====================================================================
\section{Phụ lục Phương pháp: Kiểm soát Chất lượng (QA), KPI \& Tái lập}

\subsection{Tác động QA lên giá \& Nguyên tắc xử lý (Bài 10)}
\begin{itemize}
  \item \textbf{Tác động QA lên giá tại mốc cuối (\texttt{2026-09-16}):} Ghi nhận \textbf{114} chỗ ở thiếu giá (giữ nguyên \texttt{NA}, không tự điền). Giá trung vị trước lọc ngoại lai là \textbf{211.09 CAD/đêm} (trung bình \textbf{289.01 CAD/đêm}); sau khi gắn cờ ngoại lai theo quy tắc IQR tính riêng từng \texttt{room\_type} kết hợp trần \texttt{qa.price\_max = 2,000 CAD}, giá trung vị hợp lệ đạt \textbf{201.07 CAD/đêm} ($n = 783$).
  \item \textbf{Nguyên tắc QA:} Tách riêng giá trị thiếu (\texttt{missing\_price}, \texttt{missing\_rating}, \texttt{missing\_minimum\_nights}) khỏi lỗi ngoài miền (\texttt{invalid\_price}, \texttt{invalid\_rating} ngoài \texttt{[0, 5]}, \texttt{invalid\_minimum\_nights}); chỉ gắn cờ (\texttt{flag\_only}) thay vì xóa dòng để bảo toàn cỡ mẫu cho K2--K8; phân biệt tổng lượt gắn cờ với số hàng cần kiểm tra (\texttt{rows\_with\_any\_flag}).
\end{itemize}

\begin{table}[H]
  \centering
  \small
  \renewcommand{\arraystretch}{1.2}
  \begin{tabularx}{0.76\textwidth}{Xr}
    \toprule
    \rowcolor{NavyDark}
    \color{white}\textbf{Quy tắc QA (Top 8 theo tổng 4 mốc)} &
    \color{white}\textbf{Tổng lượt/hàng bị gắn cờ} \\
    \midrule
    \texttt{review\_id\_repeated} & 29,295 \\
    \rowcolor{LightBg}
    \texttt{rows\_with\_any\_flag} & 4,128 \\
    \texttt{comment\_markup} & 3,650 \\
    \rowcolor{LightBg}
    \texttt{missing\_price} & 213 \\
    \texttt{price\_outlier} & 169 \\
    \rowcolor{LightBg}
    \texttt{missing\_rating} & 121 \\
    \texttt{room\_change} & 9 \\
    \rowcolor{LightBg}
    \texttt{missing\_minimum\_nights} & 1 \\
    \bottomrule
  \end{tabularx}
  \caption{\small Tổng hợp Top 8 quy tắc kiểm soát chất lượng (QA) trên cả 4 mốc chụp.}
\end{table}

\subsection{Định nghĩa 8 Chỉ số KPI (K1--K8)}

\begin{table}[H]
  \centering
  \footnotesize
  \renewcommand{\arraystretch}{1.22}
  \begin{tabularx}{\textwidth}{l >{\raggedright\arraybackslash}p{4.3cm} >{\raggedright\arraybackslash}p{2.3cm} >{\raggedright\arraybackslash}X}
    \toprule
    \rowcolor{NavyDark}
    \color{white}\textbf{Mã} &
    \color{white}\textbf{Tên chỉ số} &
    \color{white}\textbf{Đơn vị} &
    \color{white}\textbf{Công thức / Điều kiện chính} \\
    \midrule
    \textbf{K1} & Giá đêm trung vị & CAD/đêm & \texttt{median(price} hợp lệ\texttt{)} \\
    \rowcolor{LightBg}
    \textbf{K2} & Tỷ lệ ngày không khả dụng theo vùng/mốc & \% & \texttt{100*(1-mean(available))} \\
    \textbf{K3} & Tỷ lệ ngày không khả dụng theo tháng lịch & \% & \texttt{100*(1-mean(available))} theo tháng ngày lịch \\
    \rowcolor{LightBg}
    \textbf{K4} & Tỷ lệ yêu cầu thuê dài hạn & \% & \texttt{100*mean(minimum\_nights>=30)} \\
    \textbf{K5} & Tốc độ review trong mẫu & sample reviews / listing / quarter & review duy nhất/quý / listing duy nhất trong mẫu \\
    \rowcolor{LightBg}
    \textbf{K6} & Listing chưa có review & \% & \texttt{100*mean(number\_of\_reviews==0)} \\
    \textbf{K7} & Cảm xúc tiêu cực và điểm khía cạnh & \% hoặc \texttt{[-1,1]} & \texttt{negative/valid}; \texttt{(praise-complaint)/(praise+complaint)} \\
    \rowcolor{LightBg}
    \textbf{K8} & Listing khả dụng cao có nhiều review & \% & \texttt{mean(availability>=threshold and reviews>=threshold)*100} \\
    \bottomrule
  \end{tabularx}
  \caption{\small Định nghĩa chuẩn hóa 8 chỉ số KPI sử dụng trong báo cáo.}
\end{table}

\subsection{Lệnh Tái lập \& Trạng thái Bàn giao}
\begin{samepage}
\begin{storybox}[Tái lập \& Kiểm tra tự động]
\begin{itemize}
  \item \textbf{Chạy toàn bộ pipeline 1 lệnh:} \texttt{python run\_pipeline.py --config configs/Toronto.yaml} (thêm \texttt{--skip-download --skip-llm} khi kiểm thử ngoại tuyến).
  \item \textbf{Kiểm tra tự động:} \texttt{pytest} (35 bài test), \texttt{check\_config\_swap.py}, \texttt{check\_report\_numbers.py} (đối chiếu tự động mọi con số trong \texttt{report.md} với \texttt{summary\_numbers.csv}), và \texttt{check\_security.py}.
  \item \textbf{Provenance bản nháp:} mã nền trước khi sửa báo cáo là \texttt{2215a69df40f3c83b09197ce8e8ff26d21681209}; \texttt{working\_tree\_dirty = True}. Trước khi nộp, nhóm cần commit bản cuối, cập nhật lại mã này và gắn tag \texttt{final}. Pipeline không tự động \texttt{commit}, \texttt{tag} hay \texttt{push}.
\end{itemize}
\end{storybox}
\end{samepage}

\end{document}
